\documentclass{article}
\usepackage{mathtools} 
\usepackage{fontspec}
\usepackage[UTF8]{ctex}
\usepackage{amsthm}
\usepackage{mdframed}
\usepackage{xcolor}
\usepackage{amssymb}
\usepackage{amsmath}


% 定义新的带灰色背景的说明环境 zremark
\newmdtheoremenv[
  backgroundcolor=gray!10,
  % 边框与背景一致，边框线会消失
  linecolor=gray!10
]{zremark}{说明}

% 通用矩阵命令: \flexmatrix{矩阵名}{元素符号}{行数}{列数}
\newcommand{\flexmatrix}[4]{
  \[
  #1 = \begin{pmatrix}
    #2_{11}     & #2_{12}     & \cdots & #2_{1#4}   \\
    #2_{21}     & #2_{22}     & \cdots & #2_{2#4}   \\
    \vdots      & \vdots      & \ddots & \vdots     \\
    #2_{#31}    & #2_{#32}    & \cdots & #2_{#3#4}
  \end{pmatrix}
  \]
}

% 简化版命令（默认矩阵名为A，元素符号为a）: \quickmatrix{行数}{列数}
\newcommand{\quickmatrix}[2]{\flexmatrix{A}{a}{#1}{#2}}

\begin{document}
\title{习题 16.3}
\author{张志聪}
\maketitle

\section*{1}

\begin{itemize}
  \item (1)

        令$\varphi(x, y) = x^2 + y^2$，
        $\varphi(x, y)$在$\mathbb{R}^2$连续；

        $tan x$在定义域$D \subset \mathbb{R}$上连续，
        由复合函数的连续性（更确切的说，是定理16.7的变形，但证明方法是一致的）可知，
        $tan(\varphi(x, y))$连续当且仅当$\varphi(x, y) \in D$。
\end{itemize}

\section*{3}

说明：相比于定理16.6来说，无需知道$\lim\limits_{y \to b} \lim\limits_{x \to a} f(x)$存在；

证明：

与命题16.6的证明类似。

设$P_0 = (a, b)$，
$\lim\limits_{(x, y) \to (a, b)} = A$，由定义可知，对任意$\epsilon> 0$，存在$\delta > 0$，
使得当$P(x, y) \in U^{\circ}(P_0; \delta)$时，有
\begin{align*}
  |f(x, y) - A| < \epsilon
\end{align*}

不妨设$U^{\circ}(P_0; \delta)$在$y$轴上的投影$Y$属于$2^\circ$中邻域$B$，
即$Y \subset B$，
故对任意$y \in Y$，都有
\begin{align*}
  0 < |y - b| < \delta
\end{align*}
已知$x \to a$，有
\begin{align*}
  \lim\limits_{x \to a} f(x, y) = \varphi(y)
\end{align*}
故
\begin{align*}
  |\varphi(y) - A| < \epsilon
\end{align*}
注：可以理解为$x$足够靠近$a$.

由$\epsilon$的任意性和$\delta$的存在性可知
\begin{align*}
  \lim\limits_{y \to b} \lim\limits_{x \to a} f(x)
\end{align*}

\end{document}